% Error sources in automated glacier delineation and how the pipeline handles them.
% Standalone: compiles on its own (e.g. Overleaf). To use in a paper, copy the preamble
% packages and the sidewaystable environment.
\documentclass[10pt]{article}
\usepackage[margin=1.8cm]{geometry}
\usepackage{booktabs}
\usepackage{tabularx}
\usepackage{threeparttable}
\usepackage{rotating}
\usepackage{array}
\usepackage{wasysym}

\newcolumntype{L}[1]{>{\raggedright\arraybackslash}p{#1}}
\newcolumntype{Y}{>{\raggedright\arraybackslash}X}
\newcommand{\stfull}{\CIRCLE}
\newcommand{\stpart}{\LEFTcircle}
\newcommand{\stnone}{\Circle}
\newcommand{\grp}[1]{\multicolumn{6}{l}{\textit{#1}}\\}

\begin{document}

\begin{sidewaystable}
\centering
\footnotesize
\setlength{\tabcolsep}{4pt}
\renewcommand{\arraystretch}{1.18}
\begin{threeparttable}
\caption{Error sources in multi-temporal glacier delineation of Chhota Shigri Glacier, the
pipeline stage that addresses each, its implementation, and the residual effect on the
products. Satellite images are never edited; each source is handled by data rejection,
model classification, temporal aggregation or post-processing.}
\label{tab:error_sources}
\begin{tabularx}{\textheight}{L{2.5cm} L{3.3cm} L{3.2cm} Y L{4.6cm} c}
\toprule
\textbf{Error source} & \textbf{Effect on delineation} & \textbf{Mitigation (stage)} &
\textbf{Implementation and parameters} & \textbf{Residual effect (quantified)} & \textbf{Status}\tnote{a} \\
\midrule

\grp{Atmospheric}
Cloud &
Occludes the surface; bright cloud is confused with snow and ice &
Scene rejection and per-pixel masking (acquisition; inference) &
s2cloudless, $p \ge 0.4$, 4\,px averaging, 2\,px (40\,m) dilation. Scene rejected if domain cloud
fraction $> 0.5$ (SCL pre-screen $> 0.6$). Cloudy pixels set to NaN after prediction and excluded
from the per-pixel median &
Pixels with no clear observation are gap-filled (see \emph{Data gaps}) &
\stfull \\
Cloud shadow &
Darkens ice; may be mapped as non-glacier &
Indirect (acquisition; aggregation) &
40\,m cloud-mask dilation covers adjacent shadow; per-pixel median and scene-outlier rejection.
SCL shadow class (3) is computed but not used as a per-pixel mask in the final ensemble &
Shadow distant from its cloud can enter the median; not quantified &
\stpart \\
\addlinespace

\grp{Surface conditions}
Seasonal and fresh snow &
Snow cover outside the glacier inflates or merges the outline; autumn snowfall hides margins &
Acquisition window and snow screening (acquisition) &
Core window 15 Aug--15 Sep in every year. Scene rejected if SCL snow covers $> 45\%$ of the
domain outside the RGI7 outline (fresh-snowfall indicator) &
Scenes after 15 Sep mapped 1.5--4.6\,km$^2$ less glacier; the 15 Aug--15 Oct window lowers
inventory IoU from 0.816 to 0.778 and adds a spurious $-2.65$\,km$^2$ step (2023--2025) &
\stfull \\
Thin snow on debris margins &
Margins alternate between snow-covered and exposed debris between scenes &
Multi-date ensemble (aggregation); exclusion from training labels (self-training) &
Per-pixel median of calibrated $P(\mathrm{glacier})$ over all clear scenes of a year. Pseudo-labels
ignore all temporally unstable pixels (34.7\% of the domain) &
Single-scene area spread $\approx 1$\,km$^2$, reduced to 0.3--0.8\,km$^2$ per epoch; 1.74\,km$^2$
of unstable margin pixels vs 0.94\,km$^2$ persistent loss. Dominant limit on area-change detection &
\stpart \\
Debris-covered ice &
Spectrally similar to surrounding rock; terminus and margins are ambiguous &
Model inputs and regional fine-tuning (model) &
DEM channels (elevation, slope); transformer global context; HMA fine-tuning on debris-rich
tiles &
DEM adds $+0.149$ IoU [0.078, 0.241]. Absolute terminus position differs by up to $\sim$800\,m
between models; only terminus \emph{change} is reported &
\stpart \\
Bare rock and moraine &
False-positive glacier on bright rock; false negatives at ice--rock contacts &
Model classification; object filtering (model; post-processing) &
Six optical bands (incl.\ SWIR B11, B12) and DEM; objects $< 100$\,px (0.01\,km$^2$) removed,
holes $< 50$\,px filled &
Included in the HMA held-out accuracy (M2 IoU 0.779) &
\stfull \\
Surface water (proglacial lakes) &
Water may be confused with shadowed ice or ice &
None &
No water mask is applied; separation relies on the model &
Not quantified &
\stnone \\
\addlinespace

\grp{Illumination and timing}
Terrain shadow &
Steep shaded slopes darken ice and rock &
Indirect (acquisition; model; aggregation) &
Window closes before low-sun conditions (after mid-September); slope input; median and
scene-outlier rejection. No explicit terrain-shadow mask &
Not separately quantified &
\stpart \\
Seasonal timing between epochs &
Different stages of the melt season produce apparent change &
Fixed acquisition window (acquisition) &
Identical calendar window for every epoch &
Both 2025 scenes are from August whereas 2023 includes September; part of the 2023--2025 step
may be seasonal &
\stpart \\
\addlinespace

\grp{Geometric and radiometric}
Misregistration &
Image offsets produce false advance or retreat at the outline &
Co-registration (preprocessing) &
Phase cross-correlation (upsampling 20) on Sobel gradients of B8, 256\,px windows, on stable
terrain (outside domain $+300$\,m and all glaciers $+150$\,m, slope $< 35^\circ$, NDSI $< 0.2$,
cloud-free); Fourier shift if $> 0.25$\,px &
Two scenes corrected (0.67$\to$0.11\,px; 0.36$\to$0.14\,px); all residuals $\le 1.6$\,m;
residual enters terminus $\sigma = \sqrt{\mathrm{res}^2 + \mathrm{coreg}^2}$ &
\stfull \\
Radiometric baseline change &
ESA processing baseline 04.00 (2022) shifts reflectance by $-1000$\,DN &
Harmonisation (acquisition) &
$\rho = (\mathrm{DN} + \mathrm{RADIO\_ADD\_OFFSET}) / \mathrm{QUANTIFICATION\_VALUE}$, read per
product from the L1C metadata &
None expected &
\stfull \\
\addlinespace

\grp{Temporal aggregation}
Anomalous whole scenes &
One scene (fresh snow partly flagged as cloud, low sun) biases the yearly product &
Robust scene-outlier rejection (aggregation) &
Scene rejected if $|A_i - \tilde{A}| > \max(3 \times 1.4826\,\mathrm{MAD},\ 0.5\,\mathrm{km}^2)$ &
Rejected scenes are logged per model and epoch &
\stfull \\
Data gaps &
Pixels never observed clear in the window are undefined &
Same-year gap filling (aggregation) &
Filled from clear scenes of the \emph{same} year outside the core window (no cross-epoch
leakage); remaining pixels reported as unobserved &
2017 and 2019 have 2 core-window scenes; $\approx 3$\,km$^2$ of the domain has a single clear
look &
\stpart \\
\addlinespace

\grp{Spatial context and noise}
Neighbouring glaciers &
Adjacent ice bodies merge with the target glacier &
Analysis domain and connectivity (post-processing) &
Domain = RGI7 $+1$\,km minus 33 neighbouring RGI7 glaciers (47.9\,km$^2$). Keep 8-connected
components intersecting the RGI7 core (outline eroded by 300\,m) &
Core identifies the ice body only; margins and terminus are model-derived &
\stfull \\
Classification noise &
Isolated pixels and slivers along the outline &
Morphological filtering (post-processing) &
Threshold 0.5; $3 \times 3$ median filter; slivers $< 4$\,px removed after polygonisation;
\texttt{make\_valid} and raster--vector area check &
Included in the delineation uncertainty ($\tfrac{1}{2}$\,px $\times$ perimeter) &
\stfull \\
\addlinespace

\grp{SAR (model M3 only)}
Layover and radar shadow &
No valid backscatter on steep slopes &
Value filling (acquisition) &
Missing $\sigma^0$ filled with the composite median &
SAR gave no benefit: $-0.016$ IoU [$-0.029$, $-0.008$] vs optical $+$ DEM &
\stpart \\
\bottomrule
\end{tabularx}
\begin{tablenotes}\footnotesize
\item[a] \stfull\ addressed explicitly; \stpart\ partly or indirectly addressed; \stnone\ not addressed.
\item IoU values are on held-out HMA benchmark tiles (57 tiles; 22 tiles for the SAR and DEM
ablations), with 95\% paired-bootstrap confidence intervals in brackets. Areas are multi-date
ensemble products on the 10\,m UTM 43N grid. MAD: median absolute deviation; SCL: Sentinel-2
scene classification layer; NDSI: normalised difference snow index.
\end{tablenotes}
\end{threeparttable}
\end{sidewaystable}

\end{document}
